\documentclass[UTF8,a4paper,11pt,openany]{ctexbook}
\usepackage{../common/statlearnbook}
\SLSetBookMetadata{第 5 册}{统计学习方法讲义 v2.7.0}{采样方法、主题模型与图排序}
\begin{document}
\setcounter{page}{662}
\setcounter{chapter}{33}
\setcounter{figure}{0}
\pagestyle{headings}
\chaptermark{Gibbs 抽样}
\sectionmark{本章地图与依赖路线}
\noindent 图\ref{fig:V5-C04-dependency-graph}概括从联合目标与局部因子到 Monte Carlo 诊断的六节点学习与计算主链。
% v2.7.0 figure UID: FIG-P630-01
% R2: six-node dependency chain; every segment is an independent directed edge.
\tikzset{slfig-FIG-P630-01/.style={font=\fontsize{9.6pt}{11.6pt}\selectfont,every path/.append style={line cap=round,line join=round}}}
\begin{figure}[htbp]
\centering
\begin{tikzpicture}[slfig-FIG-P630-01,>=Stealth,every node/.style={align=center},
  core/.style={draw=SLBlue,rounded corners=2.2pt,fill=SLBlue!6,
    minimum width=31mm,minimum height=13mm,inner xsep=1.5mm,inner ysep=1.2mm,
    font=\fontsize{9.6pt}{11.8pt}\selectfont,line width=.72pt},
  side/.style={draw=SLRuleGray,rounded corners=2.2pt,fill=SLSoftGray,
    text width=22.5mm,minimum height=16mm,inner xsep=1.5mm,inner ysep=1.0mm,
    font=\fontsize{9.6pt}{11.6pt}\selectfont,line width=.55pt},
  boundary/.style={draw=SLGold,rounded corners=2.2pt,fill=SLGold!8,
    minimum width=76mm,minimum height=9mm,inner xsep=2mm,inner ysep=1mm,
    font=\fontsize{10.0pt}{12.0pt}\selectfont\bfseries,line width=.72pt},
  flow/.style={draw=SLBlue,line width=.95pt,-{Stealth[length=2.0mm,width=1.35mm]}},
  leader/.style={draw=SLGray,line width=.50pt}]
  \node[core] (joint) at (-3.80,.95) {联合目标 / 局部因子};
  \node[core] (cond) at (0,.95) {给定 $x_{-j}$ 的满条件\\$\pi_j(\cdot\mid x_{-j})$};
  \node[core] (coord) at (3.80,.95) {单坐标核 $K_j$\\只更新 $x_j$};
  \node[core] (scan) at (3.80,-.95) {扫描核\\系统 / 随机};
  \node[core] (sample) at (0,-.95) {相关样本};
  \node[core] (diag) at (-3.80,-.95) {诊断\\MCSE / ESS / 轨迹};
  \draw[flow] (joint.east)--(cond.west);
  \draw[flow] (cond.east)--(coord.west);
  \draw[flow] (coord.south)--(scan.north);
  \draw[flow] (scan.west)--(sample.east);
  \draw[flow] (sample.west)--(diag.east);
  \node[side] (correct) at (-6.40,2.85) {正确性条件\\目标保持\\支持可达\\遍历性};
  \node[side] (mix) at (6.40,-2.70) {混合效率\\自相关长度\\有效样本量};
  \draw[leader] (correct.south east)--(joint.north west);
  \draw[leader] (mix.north west)--(scan.south east);
  \node[boundary] at (0,-3.10)
    {正确内核 $\ne$ 快速混合};
\end{tikzpicture}
\captionsetup{width=.94\linewidth}
\caption{满条件把联合目标转为单坐标更新，扫描后得到需以MCSE、ESS与轨迹诊断的相关样本}
\label{fig:V5-C04-dependency-graph}
\end{figure}

\FloatBarrier
\noindent\textbf{读图顺序。}先看给定其余坐标$x_{-j}$后，满条件$\pi_j(\cdot\mid x_{-j})$如何只重抽$x_j$；再看单坐标核$K_j$怎样组成系统或随机扫描核并产生相关样本；最终用MCSE、ESS与轨迹诊断评估统计误差和混合效率。主链箭头只表示学习与计算依赖，不是概率图的生成时间方向；目标保持仍不自动推出不可约、收敛或快速混合。\par
\end{document}
